\documentclass[10pt,a4paper]{amsart}
\usepackage[margin=19mm]{geometry}
\usepackage{amsmath,amssymb,mathtools}
\usepackage[scr=rsfs,cal=euler]{mathalfa}
\usepackage{xfrac}
\usepackage[unicode,hidelinks,bookmarks=false]{hyperref}
\newcommand{\ie}{i.e.\ }
\newcommand{\cf}{cf.\ }
\newcommand{\resp}{respectively\ }
\newcommand{\Subsection}{\subsection}
\providecommand{\nobreakdash}{\nobreak\hbox{-}}
\setlength{\emergencystretch}{2em}
\multlinegap0pt
\usepackage{siunitx}
\usepackage{siunitx}
\usepackage{siunitx}
\usepackage{siunitx}
\usepackage{amssymb}
\usepackage{xcolor}
\usepackage{algorithm}\usepackage{algpseudocode}
\usepackage{bm}
\usepackage{mathtools}
\usepackage{siunitx}
\newtheorem{lemma}{Lemma}
\title{A Limit Theory of Foundation Models: \\ A Perspective of Deciphering the Phenomenon of Intelligent Emergence and Scaling Laws}
\author{Jun Shu, Junxiong Jia, Deyu Meng, Zongben Xu}
\date{Source: arXiv:2604.24037v4 (CC BY 4.0)}
\begin{document}
\maketitle

\noindent\emph{Source and licence.} A Limit Theory of Foundation Models: \\ A Perspective of Deciphering the Phenomenon of Intelligent Emergence and Scaling Laws. Version arXiv:2604.24037v4 (CC BY 4.0), distributed by arXiv under the Creative Commons Attribution 4.0 International licence, http://creativecommons.org/licenses/by/4.0/, as recorded in the arXiv licence metadata for this exact version and shown on the abstract page https://arxiv.org/abs/2604.24037v4. Full licence notice: \texttt{licenses/arxiv-2604.24037-CC-BY-4.0.txt}.\par\smallskip

\noindent\textit{Editorial scope.} Counted unit: the article's lemma lemma (source0.tex lines 5222--5240). The article's complete original proof of it is delivered with it (source0.tex lines 5242--5315, one \texttt{proof} environment of 74 lines). No mathematics is written, repaired or re-derived: the statement and the proof are the article's own lines, in the article's own notation, and no retained line is substituted or re-flowed.  No further source-local context is needed: every cross-reference of every retained line resolves inside the two delivered spans.  Nothing was omitted from any retained span, so the package carries no omission record.  No retained line contains a citation, so the wrapper carries no bibliography and no bibliography entry.  No retained line contains a figure environment, an image-inclusion command or drawing code, so no image is delivered and the package declares no asset dependency.  Editorial formatting only, in addition: an AMS \texttt{amsart} preamble that restates the article's own declarations and macros used by the retained lines, namely the environments \texttt{lemma} on separate counters, together with the standard packages those lines need.  The retained environments print in this document as Lemma 1 in order of appearance, whereas the article prints the same items with the numbering of its own full document.  The complete native source of the article as distributed in the arXiv e-print for this exact version is bundled unchanged as \texttt{sources/199206/source0.tex}.
\begin{lemma}
	Let
	\(
	U=T_1^{\alpha_1}T_2^{\alpha_2}\cdots T_r^{\alpha_r},
	\)
	where \(\alpha_1,\ldots,\alpha_r\in \{0,1\}\), and set
	\(
	m=\sum_{i=1}^r \alpha_i .
	\)
	Then \(U\) is \(\frac{\gamma}{r}\)-strongly quasi-nonexpansive, that is, for every
	\(z \in \bigcap_{i=1}^r \mathrm{Fix}(T_i)\) and every \(x\in \bigcap_{i=1}^r \mathbb{D}(T_i)\),
	\[
	\|Ux-z\|^2
	\leq
	\|x-z\|^2
	-
	\frac{\gamma}{r}\|Ux-x\|^2 .
	\]
\end{lemma}

\begin{proof}
	We write \(U\) as a composition of \(m\) mappings:
	\[
	U=S_rS_{r-1}\cdots S_1,
	\]
	where each \(S_j\) is equal to one of the mappings \(T_i\) or $I$. For a fixed
	\(x\in \bigcap_{i=1}^r \mathbb{D}(T_i)\), define
	\[
	x_0=x,\qquad x_j=S_jx_{j-1},\quad j=1,\ldots,m.
	\]
	Then \(x_r=Ux\). Since \(z\in F\), we have \(z\in \operatorname{Fix}(S_j)\)
	for every \(j=1,\ldots,r\). Hence, by the \(\gamma\)-strong
	quasi-nonexpansiveness of \(S_j\), we obtain
	\[
	\|x_j-z\|^2
	=
	\|S_jx_{j-1}-z\|^2
	\leq
	\|x_{j-1}-z\|^2
	-
	\gamma\|S_jx_{j-1}-x_{j-1}\|^2 .
	\]
	Equivalently,
	\[
	\|x_j-z\|^2
	\leq
	\|x_{j-1}-z\|^2
	-
	\gamma\|x_j-x_{j-1}\|^2 .
	\]
	Summing this inequality over \(j=1,\ldots,r\), we get
	\[
	\|x_m-z\|^2
	\leq
	\|x_0-z\|^2
	-
	\gamma\sum_{j=1}^r \|x_j-x_{j-1}\|^2 .
	\]
	Since \(x_0=x\) and \(x_r=Ux\), this gives
	\[
	\|Ux-z\|^2
	\leq
	\|x-z\|^2
	-
	\gamma\sum_{j=1}^r \|x_j-x_{j-1}\|^2 .
	\]
	On the other hand,
	\[
	Ux-x=x_r-x_0
	=
	\sum_{j=1}^r (x_j-x_{j-1}).
	\]
	By the Cauchy--Schwarz inequality,
	\[
	\|Ux-x\|^2
	\leq
	r\sum_{j=1}^r \|x_j-x_{j-1}\|^2 .
	\]
	Thus,
	\[
	\sum_{j=1}^r \|x_j-x_{j-1}\|^2
	\geq
	\frac{1}{r}\|Ux-x\|^2 .
	\]
	Substituting this estimate into the previous inequality yields
	\[
	\|Ux-z\|^2
	\leq
	\|x-z\|^2
	-
	\frac{\gamma}{r}\|Ux-x\|^2 .
	\]
	Therefore, \(U\) is \(\frac{\gamma}{r}\)-strongly quasi-nonexpansive.
\end{proof}

\end{document}
